\section{Realization of maximal closure claims}\label{sec:realization}
This section corresponds to the application of M\'emin's closedness theorem after component convergence in the proof of \cite[Theorem 4.2]{DS}. We use \contractref{C1} to realize the same limit and its terminal value as an original-price integral.

The preceding construction yields gains $R_n=\widetilde M_n+\widetilde A_n$
in the original market, converging almost surely uniformly over all time to
$X\geq-1$ with $X_\infty=h$. Their martingale components are elementary-Émery
Cauchy; the variations of their finite-variation differences are Cauchy in
probability on every finite horizon.
\begin{theorem}[Maximal-claim realization]\label{prop:realization}
Under NFLVR every almost-sure maximal $h\in\clprob{K_1}$ belongs to $K_1$.
\end{theorem}
\begin{proof}
Apply the test-uniform M expectation Cauchy estimate for the same convexified sequence with error $\varepsilon^2$.
Markov's inequality gives $Q_*(e_T^J(\widetilde M_m,\widetilde M_n)>\varepsilon)\leq\varepsilon$ with one cutoff for all tests.
The preceding maximality contradiction supplies the variation probability estimate for the FV differences.
The same $R_n$ are strongly adapted and c\`adl\`ag, and differences of the $\widetilde A_n$ are locally BV.
Together with the retained all-time uniform limit $X\geq-1$ and its own terminal $h$, these verify every hypothesis of \contractref{C1}.
Hence $h\in K_1$.
\end{proof}
